\documentclass[11pt]{article}
\usepackage[utf8]{inputenc}
\usepackage[T1]{fontenc}
\usepackage{geometry}
\geometry{a4paper, margin=1in}
\usepackage{booktabs}
\usepackage{amsmath}
\usepackage{array}
\usepackage{hyperref}
\hypersetup{colorlinks=true, linkcolor=blue, citecolor=blue, urlcolor=blue}

\title{Predictive Model and Proof-of-Concept Plan:\\
Forecasting Flight Delays and Validating on the 2025 Government Shutdown}
\author{Pranish Uprety \\ Advisor: Dr. Sushil Poudel}
\date{\today}

\begin{document}
\maketitle

\section{Goal}
The objective is a publishable proof of concept: a flight-delay prediction model
built on a set of meaningful input factors, validated by predicting a known
real-world disruption (the 2025 federal government shutdown) and comparing the
model's output against what actually happened. The lower the prediction error on
this held-out disruption, the stronger the evidence that the model captures the
real drivers of delay. We focus on the New York metropolitan sub-network
(EWR, JFK, LGA) within the Northeast corridor, where we have already shown that
delay propagates directionally (Newark drives JFK and LaGuardia roughly one hour
later, Pearson $r \approx 0.57$, BTS January 2023).

\section{The Predictive Model}
\paragraph{Target.} The model predicts an airport's delay for a future time
window (next hour, and later next day), measured as mean departure delay in
minutes per airport-hour.

\paragraph{Inputs (key factors).} The model combines local, network, weather,
operational, staffing, and economic factors (Table~\ref{tab:factors}).

\paragraph{Model progression.} We start simple and add complexity only if it
helps: (1) regularized regression as a baseline; (2) gradient-boosted trees for
nonlinear interactions; (3) the spatio-temporal graph neural network with the
intervention head, which is the eventual contribution. Each step is compared on
the same held-out data so improvements are attributable.

\begin{table}[h]
\centering
\small
\caption{Key factors (features) for the delay-prediction model.}
\label{tab:factors}
\begin{tabular}{>{\raggedright}p{2.6cm} >{\raggedright}p{6.2cm} p{3.4cm}}
\toprule
\textbf{Group} & \textbf{Factors} & \textbf{Source} \\
\midrule
Temporal & hour of day, day of week, holiday, season & BTS schedule \\
Local history & airport's own delay at $t{-}1,t{-}2,t{-}3$ & BTS on-time \\
Network (graph) & upstream-neighbor delays (the propagation signal) & BTS on-time \\
Traffic / congestion & scheduled departures and arrivals per hour, live density & BTS, OpenSky \\
Weather & wind, visibility, ceiling, precipitation; or BTS WeatherDelay & NOAA METAR / BTS \\
Operational & late-aircraft delay, carrier delay, NAS delay, taxi-out & BTS on-time \\
Staffing / disruption & ground stops, ground delay programs, shutdown indicator & FAA, news/policy \\
Economic (finance) & cost-of-delay per minute, passengers affected, \$ impact & FAA, BTS T-100 \\
\bottomrule
\end{tabular}
\end{table}

\section{Proof of Concept: the 2025 Shutdown Test}
\paragraph{Idea.} Train the model on normal operations, then ask it to predict
delays during the 2025 shutdown window, a period it was not trained on, and
compare its predictions against the actual recorded delays. Good predictions on
an unusual, high-stress period are strong evidence the model generalizes.

\paragraph{Protocol.}
\begin{enumerate}
  \item Train on pre-shutdown months (normal operations).
  \item Predict delays across the 43-day shutdown period (Oct.\ 1 onward, 2025).
  \item Compare predicted vs.\ actual BTS delays; report error.
  \item Cross-check the implied excess delay against the published Richmond Fed
        difference-in-differences estimate ($\sim$1M extra delay work-hours,
        $\sim\$66$M), an external counterfactual benchmark.
\end{enumerate}

\paragraph{Why the shutdown is ideal.} It is a large, well-documented, real
disruption with an external economic estimate to validate against, and it
directly exercises the staffing factor (controllers working unpaid, increased
sick leave, FAA-imposed flight cuts), which a purely weather-based model would
miss.

\section{Evaluation Metrics}
\begin{itemize}
  \item \textbf{MAE / RMSE} (minutes): average prediction error. Lower is better.
  \item \textbf{$R^2$}: fraction of delay variation explained on held-out data.
  \item \textbf{Feature ablations}: remove one factor group at a time to show
        which factors matter (e.g., does the network signal add value over local
        plus weather?).
  \item \textbf{Economic error}: convert predicted vs.\ actual delay into dollars
        and compare to the published shutdown cost estimate.
\end{itemize}

\section{Data To Collect (for the actual experiments later)}
\begin{table}[h]
\centering
\small
\caption{Datasets required to run the full proof of concept.}
\begin{tabular}{>{\raggedright}p{3.6cm} >{\raggedright}p{5.2cm} p{3.4cm}}
\toprule
\textbf{Data} & \textbf{Use} & \textbf{Source} \\
\midrule
BTS on-time (2022--2025) & delays, schedules, delay causes & transtats.bts.gov \\
2025 shutdown window & test period actuals & BTS \\
NOAA METAR & weather controls & weather.gov / NOAA \\
OpenSky ADS-B & live and historical traffic density & opensky-network.org \\
FAA cost-of-delay & dollar conversion & faa.gov \\
Richmond Fed brief & external counterfactual benchmark & richmondfed.org \\
\bottomrule
\end{tabular}
\end{table}

\section{This Week and Next Steps}
This week was a data-and-design week: scoping the factor list above, identifying
the data sources, and defining the shutdown proof-of-concept protocol. Next steps
are to collect the multi-year BTS slices and the 2025 shutdown window, add weather
as the first new factor, and run the train-on-normal / test-on-shutdown
comparison to get the first end-to-end error number.

\end{document}
